\section*{अध्याय \ref{ch:b3:endocrinology} --- अंतःस्रावी-विज्ञान और समस्थापन}

\begin{solution}{exo:b3:endocrinology:1}
\omterm{def:b3:endocrinology:glucose}{इंसुलिन}: पेप्टाइड, पृष्ठीय ग्राही (कोई टाइरोसिन काइनेज़), मिनटों में
काम, अर्ध-आयु काल लगभग \qty{5}{min}। कॉर्टिसोल: स्टेरॉइड, \omterm{def:b3:endocrinology:hormone}{केंद्रकीय
ग्राही}, घंटे, अर्ध-आयु काल \qty{80}{min} (अधिकांशतः किसी वाहक से बँधा)।
एड्रिनैलीन: ऐमीन, पृष्ठीय G-प्रोटीन-युग्मित ग्राही, सेकंड,
अर्ध-आयु काल लगभग \qty{2}{min}। थायरॉक्सिन: आयोडीनयुक्त ऐमीनो अम्ल,
\omterm{def:b3:endocrinology:hormone}{केंद्रकीय ग्राही}, दिन, अर्ध-आयु काल \qty{7}{d} (वाहकों से बँधा)।
वैसोप्रेसिन: पेप्टाइड, पृष्ठीय G-प्रोटीन-युग्मित ग्राही (\omterm{def:b3:renal-osmoregulation:nephron}{संग्रह नलिका}
में cAMP), मिनट, अर्ध-आयु काल लगभग \qty{15}{min}।
\end{solution}

\begin{solution}{exo:b3:endocrinology:2}
\omterm{def:b3:endocrinology:axes}{अधश्चेतक} (TRH) $\to$ \omterm{def:b3:endocrinology:axes}{पीयूष} (TSH) $\to$ अवटु (T$_{4}$), जिसमें
T$_{4}$ ऊपर की दोनों तहों का निरोध करता है। नष्ट अवटु: T$_{4}$ कम,
TSH ऊँचा। TSH स्रवित करता \omterm{def:b3:cancer-biology:hallmarks}{अर्बुद}: TSH ऊँचा और T$_{4}$ ऊँचा (वह
पुनर्भरण उस \omterm{def:b3:cancer-biology:hallmarks}{अर्बुद} को दबा नहीं पाता)। थायरॉक्सिन की अधिक गोलियाँ:
T$_{4}$ ऊँचा, TSH दबा। आयोडीन की कमी: T$_{4}$ कम या
कम-सामान्य, TSH ऊँचा, और वह ग्रंथि उसके तहत बढ़कर घेंघा बन जाती है।
\end{solution}

\begin{solution}{exo:b3:endocrinology:3}
मस्तिष्क कम ग्लूकोस से मिनटों में मर जाता है और ऊँचे से केवल वर्षों में
हानि पाता है; और विकास को अकाल भोजों से कहीं अधिक बार मिला। इसलिए किसी
गिरावट के विरुद्ध बचाव अनावश्यक-बहुल है (\omterm{def:b3:endocrinology:glucose}{ग्लूकैगॉन}, एड्रिनैलीन,
कॉर्टिसोल, वृद्धि \omterm{def:b3:endocrinology:hormone}{हॉर्मोन}) और किसी बढ़त के विरुद्ध अकेला।
पूर्वानुमान: बहुत अधिक \omterm{def:b3:endocrinology:glucose}{इंसुलिन} तीव्र रूप से घातक है (अल्पग्लूकोजी
मूर्च्छा), और बहुत कम कोई चिरकारी रोग --- और चिकित्सालय यही देखता है।
\end{solution}

\begin{solution}{exo:b3:endocrinology:4}
\omterm{def:b3:endocrinology:clock}{कालदाता} कोई ऐसा पर्यावरणी संकेत है जो किसी घड़ी को साध देता है: प्रकाश,
भोजन का समय, तापमान, व्यायाम, सामाजिक कार्यक्रम। स्वामी घड़ी प्रकाश काम
में लेती है, \omterm{def:b3:sensory-systems:retina}{दृष्टिपटल} की मेलेनॉप्सिन \omterm{def:b3:sensory-systems:retina}{गंडिका कोशिकाओं} के ज़रिए। आठ
कालक्षेत्रों के बाद वह घड़ी दिन में केवल घंटे भर खिसकती है, इसलिए
हफ़्ते भर नींद, कॉर्टिसोल, तापमान और भूख पुराने समय पर चलते हैं जबकि
परिधीय घड़ियाँ, जिन्हें भोजन फिर सेट कर देता है, अपनी ही चाल से बहती
हैं: यानी कालक्षेत्र-श्रांति, जो पूरब की ओर अधिक बुरी है, क्योंकि घड़ी
आगे करना पीछे करने से कठिन है।
\end{solution}

\begin{solution}{exo:b3:endocrinology:5}
$250$ ग्राही अधिकृत होने पर आधी-अधिकतम अनुक्रिया: $H = K_{d}\times
250/(\num{20000} - 250) = \qty{1}{nM}\times 0.0127 = \qty{12.7}{pM}$,
यानी $K_{d}$ से अस्सी गुना नीचे। \num{2000} ग्राही के साथ: $250/1750
\times \qty{1}{nM} = \qty{143}{pM}$, यानी ग्यारह गुना कम संवेदी।
लंबा ऊँचा \omterm{def:b3:endocrinology:glucose}{इंसुलिन} अपने ग्राहियों का अधोनियमन कर देता है, इसलिए किसी दिए
असर के लिए और \omterm{def:b3:endocrinology:glucose}{इंसुलिन} चाहिए: यानी इंसुलिन-प्रतिरोध का एक घटक, जो अपने
ही ऊपर पलता है।
\end{solution}

\begin{solution}{exo:b3:endocrinology:6}
$T_{4} = 13$: $1.5\,\mathrm{e}^{0.92} = \qty{3.8}{mU/L}$; $11$:
$1.5\,\mathrm{e}^{1.84} = \qty{9.4}{mU/L}$; $9$: $1.5\,\mathrm{e}^{2.76}
= \qty{24}{mU/L}$; $25$: $1.5\,\mathrm{e}^{-4.6} = \qty{0.015}{mU/L}$।
\qty{11}{pmol/L} पर T$_{4}$ अब भी \qtyrange{10}{20}{} के भीतर है जबकि
TSH अपनी ऊपरी सीमा से दुगुने से भी अधिक (अवलक्षणी अवटु-अल्पता);
$13$ पर TSH बस परास के भीतर है; $9$ पर दोनों असामान्य; और $25$ पर
दोनों उलटी ओर असामान्य।
\end{solution}

\begin{solution}{exo:b3:endocrinology:7}
$L = s\beta/(a\gamma) = 0.2\times 0.05/(0.02\times 0.1) = 5$। स्थायी
अधिकता $g^{*} = (u/a)/(1 + L) = (2/0.02)/6 = \qty{16.7}{mg/dL}$;
पुनर्भरण के बिना $u/a = \qty{100}{mg/dL}$। \omterm{def:b3:endocrinology:glucose}{इंसुलिन} की अधिकता
$i^{*} = \beta g^{*}/
\gamma = 0.05\times 16.7/0.1 = \qty{8.3}{mU/L}$।
\end{solution}

\begin{solution}{exo:b3:endocrinology:8}
विविक्तकर ऋणात्मक है: $(a - \gamma)^{2} - 4s\beta = 0.0064 - 0.04 < 0$,
यानी दोलनी।
$\omega = \sqrt{a\gamma + s\beta - (a+\gamma)^{2}/4} =
\sqrt{0.002 + 0.01 - 0.0036} = \qty{0.092}{min^{-1}}$, आवर्तकाल $2\pi/
\omega = \qty{68}{min}$; और आयाम $\mathrm{e}$ गुना गिरता है $2/(a +
\gamma) = \qty{16.7}{min}$ में। $s$ आधा करने पर: $s\beta = 0.005$,
विविक्तकर $0.0064 - 0.02 < 0$, यानी अब भी दोलनी, $\omega = \sqrt{0.007 - 0.0036}
= \qty{0.058}{min^{-1}}$, आवर्तकाल \qty{108}{min}; क्षय-काल
अपरिवर्तित; $L = 2.5$। कोई दुर्बल पाश धीरे लौटता है और बड़ी स्थायी
त्रुटि थामे रखता है।
\end{solution}

\begin{solution}{exo:b3:endocrinology:9}
साल भर उस प्रेडनिसोलोन ने CRH और ACTH दबाए रखा; और वह अधिवृक्क वल्कुट,
जिसे कोई उद्दीपन न मिला, सिकुड़ गया। बंद करने पर कोई बाहरी स्टेरॉइड
नहीं है, \omterm{def:b3:endocrinology:axes}{अधश्चेतक} और \omterm{def:b3:endocrinology:axes}{पीयूष} को फिर चलने में हफ़्ते लगते हैं, और वह सिकुड़ी
अधिवृक्क ACTH आ भी जाए तो उस पर अनुक्रिया न दे पाती। कोई संक्रमण
सामान्य से दस गुना कॉर्टिसोल माँगता है; और कुछ भी न हो तो वाहिकाएँ
फैलती हैं, दाब तथा ग्लूकोस गिरते हैं: यानी अधिवृक्क संकट। नापने पर:
ACTH कम (वह घाव केंद्रीय है), कॉर्टिसोल कम, और अंतःक्षिप्त ACTH पर
कॉर्टिसोल की खराब बढ़त (वह ग्रंथि गल चुकी है) --- ऐडिसन रोग से उलट,
जहाँ ACTH ऊँचा होता है। इसीलिए वह धीमी कटौती।
\end{solution}

\begin{solution}{exo:b3:endocrinology:10}
एक चर: $\dot x = f(x)$, जिसमें $f' < 0$ हो, का कोई एक ही स्थिर बिंदु
$x^{*}$ है; $x > x^{*}$ के लिए $\dot x < 0$ है और $x$ एकदिष्ट रूप से
$x^{*}$ की ओर घटता है, उसे कभी पार नहीं करता (हलों की अनन्यता), और नीचे
सममित रूप से: कोई दोलन नहीं। दो चर: वह आव्यूह सम्मिश्र अभिलक्षणिक मान
रख सकता है, जैसा प्रमेय में, जब वह कारक अपने ही काल-मापक्रम वाले किसी
दूसरे चर के ज़रिए काम करे। कोई तात्क्षणिक दमन किसी स्थायी स्तर पर जम
जाता; और अनुलेखन तथा केंद्रकीय दमन के बीच का विलंब (अनुवादन, द्विलकन,
फ़ॉस्फ़ोरिलन, आयात) उस प्रोटीन को जीन के बंद होने से पहले ऊपर उछलने और
उसके फिर चालू होने से पहले नीचे उतरने देता है, और वह आवर्तकाल मोटे तौर
पर उस विलंब तथा उस प्रोटीन के अपघटन के समय के योग का दुगुना है। PER का
तेज़ अपघटन उस आवर्तकाल को छोटा कर देता है: मानव PER2 का कोई अस्थिरकारी
उत्परिवर्तन लगभग 20 घंटे की घड़ी देता है और कोई ऐसा परिवार जो शाम 7 बजे
सो जाता है।
\end{solution}

\begin{solution}{exo:b3:endocrinology:11}
चरघातांक $k(\text{Ca} - 1.2)$: $1.10$ पर $\mathrm{e}^{-2}$, PTH $= 1/
(1 + 0.135) = 0.88$; $1.15$: $0.73$; $1.20$: $0.50$; $1.25$: $1/(1 +
2.72) = 0.27$; $1.30$: $1/(1 + 7.39) = 0.12$। \omterm{thm:b3:endocrinology:loop}{नियत बिंदु} पर ढाल
$-k/4 = -5$ प्रति mmol/L: यानी हर \qty{0.01}{mmol/L} पर अधिकतम का
\qty{5}{\%}। कैल्सियम को कुछ ही प्रतिशत के भीतर थामना पड़ता है और उसका
कारक किसी विशाल बफ़र (अस्थि) पर बिना ऊपर उछले काम करता है, इसलिए कोई
तीखी अनुक्रिया सुरक्षित भी है और ज़रूरी भी। इतना तीखा कोई ग्लूकोस-पाश,
जो अपने ही सफ़ाई-काल वाले \omterm{def:b3:endocrinology:glucose}{इंसुलिन} के ज़रिए काम करे, बड़ा \omterm{thm:b3:endocrinology:loop}{पाश-लाभ} रखता
और हर भोजन के बाद गहरी अल्पग्लूकोजता में दोलन करता।
\end{solution}

\begin{solution}{exo:b3:endocrinology:12}
ग्लाइकेशन अनुत्क्रमणीय और धीमा है, इसलिए आयु $\tau$ की किसी कोशिका का
अंश $kG\tau$ ग्लाइकेटित होता है; \numrange{0}{120} दिनों में समान रूप
से बँटी आयु वाली कोशिकाओं पर औसत लेने पर HbA1c $= kG\times\qty{60}{d}$,
यानी माध्य ग्लूकोस के समानुपाती। अंशांकन: $5.5\times 60k = 0.05$ से $k =
\qty{1.5e-4}{}$ प्रति (mmol/L)(दिन); और \qty{10}{mmol/L} पर:
\qty{9.1}{\%}। 60 दिन जीती कोशिकाओं के साथ माध्य आयु 30 दिन है और उसी
ग्लूकोस के लिए HbA1c आधा रह जाता है: \qty{10}{mmol/L} वाला कोई रोगी
\qty{4.5}{\%} पढ़ेगा, यानी सामान्य।
\end{solution}

\begin{solution}{pb:b3:endocrinology:1}
\textbf{1.} \qty{90}{mg/dL} $= \qty{0.9}{g/L} = 0.9/180 = \qty{5.0}{mmol/L}$;
वितरण-आयतन में $0.9\times 15 = \qty{13.5}{g}$।
\textbf{2.} $75/15 = \qty{5}{g/L}$: यानी \qty{500}{mg/dL} की बढ़त,
\qty{28}{mmol/L}। असली शिखर इससे कहीं नीचा है क्योंकि अवशोषण घंटे भर
में फैला है जबकि निपटान ग्लूकोस को आते ही हटाता रहता है, और यकृत पहले
ही चक्कर में उसका बड़ा हिस्सा ले लेता है।
\textbf{3.} $\qty{75}{g}/\qty{60}{min} = \qty{1.25}{g/min}$;
\qty{15}{L} पर: \qty{8.3}{mg/dL/min}, यानी $u_{0}$ से चार गुना।
\textbf{4.} स्थायी अधिकता $u/a = 8.3/0.02 = \qty{417}{mg/dL}$, काल-अचर
$1/a = \qty{50}{min}$; और घंटे भर बाद उपवास से $417(1 -
\mathrm{e}^{-1.2}) = \qty{291}{mg/dL}$ ऊपर: यानी लगभग
\qty{380}{mg/dL}, \qty{21}{mmol/L}।
\textbf{5.} $L = 0.2\times 0.05/(0.02\times 0.1) = 5$; स्थायी अधिकता
$417/6 = \qty{69}{mg/dL}$, यानी \omterm{def:b3:endocrinology:glucose}{इंसुलिन} के बिना से छह गुना कम।
\textbf{6.} मस्तिष्क दिन भर में \qty{120}{g} ग्लूकोस जलाता है, कुछ भी
संचित नहीं करता और और कुछ शायद ही काम में ले सकता है: कम ग्लूकोस
मिनटों में भ्रम तथा मूर्च्छा देता है। ऊँचा ग्लूकोस प्रोटीनों का
ग्लाइकेशन करता है और वर्षों में \omterm{def:b3:sensory-systems:retina}{दृष्टिपटल}, गुर्दे तथा तंत्रिकाओं की
छोटी वाहिकाएँ और बड़ी धमनियाँ बिगाड़ देता है।
\textbf{7.} $\lambda^{2} + (a + \gamma)\lambda + (a\gamma + s\beta) = 0$:
अनुरेख $-0.12$, सारणिक $0.002 + 0.01 = 0.012$, विविक्तकर
$0.0144 - 0.048 = -0.0336$।
\textbf{8.} $\omega = \sqrt{0.0336}/2 = \qty{0.092}{min^{-1}}$; आवर्तकाल
\qty{68}{min}; क्षय-काल $2/0.12 = \qty{16.7}{min}$।
\textbf{9.} $\dot g(0) = 100(-0.06 + c\,\omega) = -2$, इसलिए $c\,\omega =
0.04$, $c = 0.44$।
\textbf{10.} $g = 0$ तब जब $\tan\omega t = -1/c = -2.29$, $\omega t =
\pi - 1.16 = 1.98$, $t = \qty{21.6}{min}$। पहले निम्निष्ठ पर, लगभग
\qty{32}{min} के पास: $100\,\mathrm{e}^{-1.94}(\cos 2.97 + 0.44\sin 2.97) =
14.4\times(-0.91) \approx -\qty{13}{mg/dL}$, यानी ग्लूकोस
\qty{77}{mg/dL}।
\textbf{11.} $t = 0$ पर $\dot i = 100\beta > 0$: \omterm{def:b3:endocrinology:glucose}{इंसुलिन} तब तक चढ़ता है
जब तक $g > 0$ हो, और तब शिखर छूता है जब $\beta g = \gamma i$ हो, जिसके
लिए $g$ का अब भी धनात्मक होना ज़रूरी है --- इसलिए वह शिखर ग्लूकोस के
उपवास तक पहुँचने से पहले आता है, और ग्लूकोस के गिरना शुरू करने के बाद
(वह $t = 0$ से गिरता है)। उस प्रतिरूप में वह शिखर \qty{12}{min} के पास
है।
\textbf{12.} घंटे भर में फैले निवेश के साथ वह शिखर तब आता है जब अवशोषण
और निपटान संतुलित हों, यानी नीचा और देर से (30--60 मिनट); वह वापसी
अवशोषण के ख़त्म होने की प्रतीक्षा करती है, इसलिए दो घंटे; और किसी
क्रमिक निवेश के लिए \omterm{def:b3:endocrinology:glucose}{इंसुलिन} का अत्युछाल छोटा होता है, इसलिए वह ढलान
हलकी।
\textbf{13.} $L = 2.5$। स्वस्थ $g^{*} = (2/0.02)/6 = \qty{16.7}{mg/dL}$;
प्रतिरोधी $100/3.5 = \qty{28.6}{mg/dL}$: यानी लगभग
\qty{12}{mg/dL} (\qty{0.7}{mmol/L}) की बढ़त, और उपवासी ग्लूकोस
\qty{102}{mg/dL}, \qty{5.7}{mmol/L} के पास --- यानी बिगड़े उपवासी
ग्लूकोस की देहली।
\textbf{14.} $i^{*} = \beta g^{*}/\gamma$: स्वस्थ $8.3$, प्रतिरोधी
\qty{14.3}{mU/L}, अनुपात $1.7$: यानी भरपाई किया गया
इंसुलिन-प्रतिरोध, लगभग-सामान्य ग्लूकोस के साथ अतिइंसुलिनता।
\textbf{15.} $L = 0.1\times 0.01/0.002 = 0.5$; $g^{*} = 100/1.5 =
\qty{66.7}{mg/dL}$, यानी स्वस्थ अवस्था से \qty{50}{mg/dL} ऊपर:
\qty{2.8}{mmol/L}, और उपवासी ग्लूकोस \qty{7.8}{mmol/L} (\qty{140}{mg/dL}),
यानी मधुमेही देहली से ऊपर।
\textbf{16.} विविक्तकर $(a - \gamma)^{2} - 4s\beta = 0.0064 - 0.004 =
0.0024 > 0$: एकदिष्ट वापसी। $\lambda = -0.06 \pm 0.0245$: यानी $-0.0355$
और $-0.0845$ प्रति मिनट; सबसे धीमा काल-अचर \qty{28}{min}, जबकि स्वस्थ
क्षय के लिए \qty{16.7}{min}: विफल होता पाश धीरे लौटता है और कभी नीचे
नहीं उतरता।
\textbf{17.} ऊँचे \omterm{def:b3:endocrinology:glucose}{इंसुलिन} के साथ ऊँचा ग्लूकोस: यानी आंशिक भरपाई वाला
प्रतिरोध --- प्ररूप~2; और दो-घंटे का मान \qty{13}{mmol/L}
\qty{11.1}{} से ऊपर है, इसलिए \omterm{def:b3:endocrinology:glucose}{मधुमेह}, केवल बिगड़ी सह्यता नहीं।
\textbf{18.} कोई C-पेप्टाइड न होने का अर्थ है कोई अंतर्जात \omterm{def:b3:endocrinology:glucose}{इंसुलिन}
नहीं: प्ररूप~1। \omterm{def:b3:endocrinology:glucose}{इंसुलिन} के बिना पेशी और वसा ग्लूकोस ले नहीं सकतीं, वसा
टूटकर वसा अम्ल तथा कीटोन बनती है, पेशी-प्रोटीन नवग्लूकोजनन के लिए
अपचयित होता है, और ग्लूकोस अपनी कैलोरियों समेत मूत्र में निकल जाता है:
वह रोगी भरपूरी के बीच भूखा मरता है।
\textbf{19.} $5\times 9/5.5 = \qty{8.2}{\%}$।
\textbf{20.} $1.5\,\mathrm{e}^{0.46\times 3} = 1.5\times 3.97 =
\qty{6.0}{mU/L}$। T$_{4}$ परास के भीतर, TSH उससे ऊपर: यानी अवलक्षणी
अवटु-अल्पता, जिसमें वह ग्रंथि विफल हो रही है और \omterm{def:b3:endocrinology:axes}{पीयूष} भरपाई कर रहा है।
\textbf{21.} $1.5\,\mathrm{e}^{0.46\times 8} = 1.5\times 39.6 =
\qty{59}{mU/L}$। कम T$_{4}$ के साथ केवल $0.3$ का TSH यह कहता है कि
\omterm{def:b3:endocrinology:axes}{पीयूष} अनुक्रिया नहीं दे रहा: वह घाव केंद्रीय है (\omterm{def:b3:endocrinology:axes}{पीयूष} या \omterm{def:b3:endocrinology:axes}{अधश्चेतक}),
अवटु में नहीं।
\textbf{22.} एक अर्ध-आयु: वह स्तर गिरकर \qty{50}{\%} रह जाता है, और वे
लक्षण हफ़्तों में सरकते हुए आते हैं। कॉर्टिसोल, जिसका अर्ध-आयु काल
\qty{80}{min} है, किसी छूटी मात्रा के घंटों के भीतर जा चुका होता है, और
उस दिन का कोई तनाव उसे बिना बचाव के मिलता है: थायरॉक्सिन का भूला हुआ
हफ़्ता असुविधाजनक है, कॉर्टिसोल का भूला हुआ दिन मार सकता है।
\textbf{23.} T$_{4}$ ऊँचा (वह \omterm{def:b3:adaptive-immunity:antibody}{प्रतिरक्षी} उस ग्रंथि को पुनर्भरण की परवाह
किए बिना चलाता है); ऊँचे T$_{4}$ से TSH दबा; और वह लघु-रैखिक संबंध उसे
अपकड़ तक ले जाता है --- $T_{4} = 30$ पर $1.5\,\mathrm{e}^{-6.9}
= \qty{0.0015}{mU/L}$ --- इसलिए कोई अपकड़ TSH ही पहला चिह्न है।
\textbf{24.} किसी \omterm{def:b3:endocrinology:hormone}{हॉर्मोन} का स्तर उसके आदेश और उसकी ग्रंथि का संतुलन
बताता है: कम T$_{4}$ के साथ ऊँचा TSH कोई विफल ग्रंथि है और ऊँचे
T$_{4}$ के साथ ऊँचा TSH कोई बेलगाम \omterm{def:b3:endocrinology:axes}{पीयूष}; ऊँचे ग्लूकोस के साथ ऊँचा
\omterm{def:b3:endocrinology:glucose}{इंसुलिन} कोई प्रतिरोधी शरीर है और ऊँचे ग्लूकोस के साथ कम \omterm{def:b3:endocrinology:glucose}{इंसुलिन} कोई
नष्ट द्वीपिका। अकेला पढ़ा जाए तो \qty{7}{pmol/L} का कोई T$_{4}$ या
\qty{30}{mU/L} का कोई \omterm{def:b3:endocrinology:glucose}{इंसुलिन} यह नहीं बताता कि दोष कहाँ है; अपने
सेनापति के साथ पढ़ा जाए तो बता देता है।
\textbf{25.} \omterm{thm:b3:endocrinology:loop}{पाश-लाभ} $L = 5$; आवर्तकाल \qty{68}{min} और क्षय-काल
\qty{17}{min}; और विफल होते पाश का उपवासी ग्लूकोस \qty{7.8}{mmol/L}।
\end{solution}
